% !TeX program = latexmk --lualatex --interaction=nonstopmode --halt-on-error --output-directory=build --synctex=1

\documentclass{article}
\usepackage[a4paper, margin=30mm]{geometry}

\usepackage[
    colorlinks,
    citecolor=black,
    linkcolor=black,
]{hyperref}
\usepackage[
  style=numeric,
  doi=true, isbn=false, url=false, eprint=false,
  date=year,% print only the year, not the month
]{biblatex}
\addbibresource{references.bib}

\usepackage{amsmath}
\usepackage{amssymb}
\numberwithin{equation}{section}

\usepackage{cancel}
\usepackage[
    per-mode=fraction
]{siunitx}


\newcommand{\dd}{\mathrm{d}}
\newcommand{\pd}[2]{\dfrac{\partial #1}{\partial #2}}
\newcommand{\pdd}[2]{\dfrac{\partial^2 #1}{\partial #2^2}}
\newcommand{\dimless}[1]{\hat{#1}}

\setlength{\parindent}{0pt}

\title{Scriven Problem in 2D and 3D}
\author{Gidon Bauer}

\begin{document}

\maketitle

\section{Original Scriven Problem}

The original Scriven problem gives the analytical solution to a spherical bubble growing in a superheated liquid due to phase-change \cite{scrivenDynamicsPhaseGrowth1995}.
This problem is three-dimensional.

\subsection{Conservation of Mass}
To derive the analytical solution, we first consider the continuity equation
\begin{equation}
    \pd{\rho}{t} + \nabla \cdot \left(\rho \vec{u}\right) = 0.
\end{equation}
Here $\rho$ is the mass density and $\vec{u}$ the velocity of the fluid.
Under the assumption of incompressibility, this reduces to
\begin{equation}
    \nabla \cdot \vec{u} = 0.
    \label{eq:continuity}
\end{equation}
Under rotational symmetry, we have $\pd{(\cdot)}{\theta} = 0$ and $\pd{(\cdot)}{\varphi} = 0$.
Then Eq. \eqref{eq:continuity} becomes
\begin{equation}
    \begin{split}
        &\frac{1}{r^2} \pd{\left(r^2 u_r\right)}{r}
        + \frac{1}{r \sin(\theta)} \cancelto{0}{\pd{\left(u_\theta \sin(\theta)\right)}{\theta}}
        + \frac{1}{r \sin(\theta)} \cancelto{0}{\pd{\left(u_\varphi\right)}{\varphi}}
        = 0 \\
        \Leftrightarrow\ &\frac{1}{r^2} \pd{\left(r^2 u_r\right)}{r} = 0 \\
        \Leftrightarrow\ &\pd{\left(u_r r^2\right)}{r} = 0. \\
        \Leftrightarrow\ &r^2 \pd{u_r}{r} + 2ru_r = 0.
    \end{split}
\end{equation}
For simplicity, we will drop the $r$ subscript of $u_r$.
We can conclude that $u r^2$ is purely a function of time, i.e. we have
\begin{equation}
    u r^2 = f(t).
    \label{eq:velocity_radius_constraint}
\end{equation}

We take the mass balance at the bubble interface \cite{scrivenDynamicsPhaseGrowth1995}
\begin{equation}
    \begin{split}
        &\pd{}{t}\left(\frac43 \pi R^3 \rho_g\right) =
        4\pi R^2 \rho_l \left[\dot{R} - u(R)\right] \\
        \Leftrightarrow\ &4\pi R^2 \rho_g \dot{R} =
        4\pi R^2 \rho_l \left[\dot{R} - u(R)\right] \\
        \Leftrightarrow\ &u(R) =
        \frac{\rho_l - \rho_g}{\rho_l} \dot{R} =
        \varepsilon \dot{R}.
    \end{split}
\end{equation}
Substituting this into Eq. \eqref{eq:velocity_radius_constraint}, we get
\begin{equation}
    u r^2 = \varepsilon \dot{R} R^2.
    \label{eq:velocity_radius_finished}
\end{equation}

\subsection{Conservation of Momentum}

The conservation of momentum for a Newtonian incompressible liquid is
\begin{equation}
    \pd{\vec{u}}{t} +
    (\vec{u} \cdot \nabla) \vec{u} +
    \frac{1}{\rho} \nabla p -
    \nu \Delta \vec{u}
    = \vec{f}.
\end{equation}
In the liquid phase, for a spherically symmetric system with no external forcing, this becomes for the radial component
\begin{equation}
    \pd{u}{t}
    + u\pd{u}{r}
    + \frac{1}{\rho_l} \pd{p}{r}
    - \nu_l \left[
        \pdd{u}{r}
        + \frac{2}{r} \pd{u}{r}
        - \frac{2u}{r^2}
    \right]
    = 0
\end{equation}

We replace the pressure $p$ by the radial component of the stress tensor
\begin{equation}
    p_{rr} = p  - 2\mu \pd{u}{r}
\end{equation}
giving us
\begin{equation}
    \pd{u}{t}
    + u\pd{u}{r}
    + \frac{1}{\rho_l} \pd{p_{rr}}{r}
    + 2 \nu_l \pdd{u}{r}
    - \nu_l \left[
        \pdd{u}{r}
        + \frac{2}{r} \pd{u}{r}
        - \frac{2u}{r^2}
    \right]
    = 0
\end{equation}


Substituting Eq. \eqref{eq:velocity_radius_finished}, we get the following expression in which the viscosity disappears
% \begin{itemize}
%     \item 
% \begin{equation}
%     \pd{u}{t} =
%     \pd{}{t} \left(\frac{\varepsilon \dot{R} R^2}{r^2}\right) =
%     \frac{\varepsilon}{r^2} \pd{}{t} \left(\dot{R} R^2\right) =
%     \frac{\varepsilon}{r^2} \left(
%         R^2 \ddot{R} +
%         2 R \dot{R}^2
%     \right)
% \end{equation}
%     \item
% \begin{equation}
%     u \pd{u}{r} =
%     \frac{\varepsilon \dot{R} R^2}{r^2}
%     \pd{}{r} \left(\frac{\varepsilon \dot{R} R^2}{r^2}\right) =
%     \frac{\varepsilon \dot{R} R^2}{r^2}
%     \left(-\frac{2 \varepsilon \dot{R} R^2}{r^3}\right) =
%     -\frac{2 \varepsilon^2 \dot{R}^2 R^4}{r^5}
% \end{equation}
%     \item
% \begin{equation}
%     \pdd{u}{r}
%     = \varepsilon \dot{R} R^2 \pdd{}{r} \frac{1}{r^2}
%     = \frac{6 \varepsilon \dot{R} R^2}{r^4}
% \end{equation}
%     \item
% \begin{equation}
%     \begin{split}
%         &\nu \left[
%             \pdd{u}{r}
%             + \frac{2}{r} \pd{u}{r}
%             - \frac{2u}{r^2}
%         \right] \\
%         =\ &\nu \left[
%             \pdd{}{r} \left(\frac{\varepsilon \dot{R} R^2}{r^2}\right)
%             + \frac{2}{r} \pd{}{r} \left(\frac{\varepsilon \dot{R} R^2}{r^2}\right)
%             - \frac{2 \varepsilon \dot{R} R^2}{r^4}
%         \right] \\
%         =\ &\nu \varepsilon \dot{R} R^2 \left[
%             \pdd{}{r} \left(\frac{1}{r^2}\right)
%             + \frac{2}{r} \pd{}{r} \left(\frac{1}{r^2}\right)
%             - \frac{2}{r^4}
%         \right] \\
%         =\ &\nu \varepsilon \dot{R} R^2 \left[
%             \frac{6}{r^4}
%             - \frac{4}{r^4}
%             - \frac{2}{r^4}
%         \right] \\
%         =\ &0
%     \end{split}
% \end{equation}
% \end{itemize}
\begin{equation}
    \frac{\varepsilon R^2 \ddot{R}}{r^2}
    + \frac{2 \varepsilon R \dot{R}^2}{r^2}
    - \frac{2 \varepsilon^2 \dot{R}^2 R^4}{r^5}
    + 2\nu_l \frac{6 \varepsilon \dot{R} R^2}{r^4}
    + \frac{1}{\rho_l} \pd{p_{rr}}{r}
    = 0
\end{equation}

We integrate this over $r$ from the bubble surface $R$ to infinity
\begin{equation}
    \begin{split}
        &\int_{R}^{\infty} \left(
            \frac{\varepsilon R^2 \ddot{R}}{r^2}
            + \frac{2 \varepsilon R \dot{R}^2}{r^2}
            - \frac{2 \varepsilon^2 \dot{R}^2 R^4}{r^5}
            + 2\nu_l \frac{6 \varepsilon \dot{R} R^2}{r^4}
            + \frac{1}{\rho_l} \pd{p_{rr}}{r}
        \right) \dd r \\
        =\ &
        \varepsilon R \ddot{R}
        + 2 \varepsilon \dot{R}^2
        - \frac{\varepsilon^2 \dot{R}^2}{2}
        + \frac{4\nu_l \varepsilon \dot{R}}{R}
        + \frac{1}{\rho_l} \left[p_{rr}\right]_R^\infty = 0 \\
        \Rightarrow\ &
        R \ddot{R}
        + \left(2 - \frac{\varepsilon}{2}\right)\dot{R}^2
        + 4\nu_l \frac{\dot{R}}{R}
        = -\frac{p_{rr}(\infty) - p_{rr}(R)}{\rho_l \varepsilon}
    \end{split}
\end{equation}

At $r \rightarrow \infty$ the radial stress is simply the ambient pressure $p_\infty$.
At the bubble interface the radial stress is
\begin{equation}
    p_{rr}(R) = p_v + p_i - \frac{2 \sigma}{R}
\end{equation}
where $\sigma$ is the surface tension,
$p_i$ the pressure of the gas phase, and
$p_v$ the partial pressure of the volatiles.

This then gives us the final equation for the conservation of momentum
\begin{equation}
    R \ddot{R}
    + \left(2 - \frac{\varepsilon}{2}\right)\dot{R}^2
    + 4\nu_l \frac{\dot{R}}{R}
    = \frac{p_v + p_i - p_\infty - \tfrac{2 \sigma}{R}}{\rho_l \varepsilon}
    \label{eq:spherically_symmetric_conservation_momentum}
\end{equation}

\subsection{Conservation of Energy}

For an incompressible fluid with negligible viscous dissipation,
the conservation of energy in temperature form can be written as
\begin{equation}
    \pd{T}{t}
    + \left(\vec{u} \cdot \nabla\right) T
    - \alpha_l \Delta T
    = \frac{Q}{\rho_l c_l}.
\end{equation}
Here $T$ is the temperature,
$\alpha_l$ the thermal diffusivity of the liquid phase,
$c_l$ the specific heat capacity of the liquid phase, and
$Q$ the volumetric heat generation.

For the spherically symmetric system, this reduces to
\begin{equation}
    \pd{T}{t}
    =
    \alpha_l \left(\pdd{T}{r} + \frac{2}{r} \pd{T}{r}\right)
    - u \pd{T}{r}
    + \frac{Q}{\rho_l c_l}
\end{equation}
and with Eq. \eqref{eq:velocity_radius_finished}, we get
\begin{equation}
    \pd{T}{t}
    =
    \alpha_l \left(\pdd{T}{r} + \frac{2}{r} \pd{T}{r}\right)
    - \frac{\varepsilon \dot{R} R^2}{r^2} \pd{T}{r}
    + \frac{Q}{\rho_l c_l}.
    \label{eq:spherically_symmetric_conservation_energy}
\end{equation}

\subsection{Mass flow}

Not used at the moment.

\subsection{Initial Conditions}

As initial conditions, we have
\begin{itemize}
    \item $R(0) = \frac{2\sigma}{p_{v,0} + p_{i,0} - p_\infty} + \delta$, where $\delta \in \mathbb{R}^+$ is a small perturbation
    \item $\dot{R}(0) = 0$
    \item $T(r, 0) = T_0$
\end{itemize}

\subsection{Boundary Conditions}

At $r \rightarrow \infty$, we have the boundary condition for the temperature
\begin{equation}
    T(\infty, t) = T_0 + \frac{1}{\rho_l c_l} \int_{0}^{t} Q(\infty, \tau) \dd \tau
\end{equation}
which reduces to $T(\infty, t) = T_0$ in case of zero heat source.

At the bubble interface, we have the boundary condition for the temperature in the one-component system
\begin{equation}
    T(R, t) = T_\text{sat}.
\end{equation}

\subsection{Simplifications}

Solving equations \eqref{eq:spherically_symmetric_conservation_momentum} and \eqref{eq:spherically_symmetric_conservation_energy} analytically is still not possible.
We will therefore introduce further simplifications.
\begin{enumerate}
    \item The viscous term in Eq. \eqref{eq:spherically_symmetric_conservation_momentum} is neglected.
    \item The inertia and surface tension term in Eq. \eqref{eq:spherically_symmetric_conservation_momentum} are neglected.
    \item The inert gas term in Eq. \eqref{eq:spherically_symmetric_conservation_momentum} is neglected.
    \item The heat source in Eq. \eqref{eq:spherically_symmetric_conservation_energy} is neglected.
\end{enumerate}
With these simplifications, we get the equations
\begin{equation}
    \frac{p_v - p_\infty}{\rho_l \varepsilon} = 0
    \quad \Rightarrow \quad
    p_v = p_\infty
\end{equation}
and
\begin{equation}
    \pd{T}{t}
    =
    \alpha_l \left(\pdd{T}{r} + \frac{2}{r} \pd{T}{r}\right)
    - \frac{\varepsilon \dot{R} R^2}{r^2} \pd{T}{r}.
    \label{eq:spherically_symmetric_conservation_energy_simplified}
\end{equation}

We write Eq. \eqref{eq:spherically_symmetric_conservation_energy_simplified} in terms of reduced temperature
\begin{equation}
    \dimless{T} = \frac{T - T_\infty}{T_\infty}.
\end{equation}
We get
\begin{equation}
    \pd{\dimless{T}}{t}
    =
    \alpha_l \left(
        \pdd{\dimless{T}}{r} +
        \frac{2}{r} \pd{\dimless{T}}{r}
    \right)
    - \frac{\varepsilon \dot{R} R^2}{r^2} \pd{\dimless{T}}{r}
    \label{eq:spherically_symmetric_conservation_energy_dimless}
\end{equation}
with initial conditions
\begin{equation}
    \dimless{T}(r, 0) = 0
\end{equation}
and boundary conditions
\begin{equation}
    \dimless{T}(R, t) = -\tau
\end{equation}
and
\begin{equation}
    \pd{\dimless{T}}{r}(R, t) = \frac{\dot{R}}{\alpha_l} \left(\xi + \omega \nu \tau\right).
\end{equation}
The dimensionless numbers are
the dimensionless superheat
\begin{equation}
    \tau = \frac{T_\infty - T_\text{sat}}{T_\infty},
    \label{eq:dimensionless_superheat}
\end{equation}
the dimensionless latent heat
\begin{equation}
    \xi = \frac{\rho_g L}{\rho_l c_l T_\infty},
    \label{eq:dimensionless_latent_heat}
\end{equation}
the dimensionless difference of heat capacity of the phases
\begin{equation}
    \nu = \frac{c_l - c_g}{c_l},
    \label{eq:dimensionless_diff_heat_capacity}
\end{equation}
and the ratio of density
\begin{equation}
    \omega = 1 - \varepsilon = \frac{\rho_g}{\rho_l}.
    \label{eq:ratio_density}
\end{equation}
Here, $L$ is the non-dimensionless latent heat of vaporization.

We then assume that $\dimless{T}(R, t) = \dimless{T}(s)$ with
\begin{equation}
    r = 2 \beta \sqrt{\alpha_l t}
    \label{eq:similarity_r}
\end{equation}
and
\begin{equation}
    s = \frac{r}{2 \sqrt{\alpha_l t}}.
    \label{eq:similarity_s}
\end{equation}

We apply this to Eq. \eqref{eq:spherically_symmetric_conservation_energy_dimless} by using
\begin{align}
    \pd{\dimless{T}}{t}
    &= \pd{\dimless{T}}{s} \pd{s}{t}
    = -\frac{r}{4 \sqrt{\alpha_l} t^{3/2}} \pd{\dimless{T}}{s}
    = -\frac{s}{2 t} \pd{\dimless{T}}{s}, \\
    \pd{\dimless{T}}{r}
    &= \pd{\dimless{T}}{s} \pd{s}{r}
    = \frac{1}{2 \sqrt{\alpha_l t}} \pd{\dimless{T}}{s}, \\
    \pdd{\dimless{T}}{r}
    &= \pd{\dimless{T}}{s} \pdd{s}{r} + \pdd{\dimless{T}}{s} \left(\pd{s}{r}\right)^2
    = \frac{1}{4 \alpha_l t} \pdd{\dimless{T}}{s}, \\
    r &= 2s \sqrt{\alpha_l t}, \\
    R &= 2 \beta \sqrt{\alpha_l t}, \\
    \dot{R} &= \beta \sqrt{\frac{\alpha_l}{t}}.
\end{align}

Then we get
\begin{equation}
    \pdd{\dimless{T}}{s}
    = -2 \pd{\dimless{T}}{s} \left(
        s
        + \frac{1}{s}
        - \frac{\varepsilon \beta^3}{s^2}
    \right)
\end{equation}

Integrating this equation over $s$ gives us
\begin{equation}
    \pd{\dimless{T}}{s} =
    \frac{A}{s^2} \exp\left(-s^2 - \frac{2 \varepsilon \beta^3}{s}\right).
\end{equation}
Applying the initial conditions and integrating again gives us
\begin{equation}
    \dimless{T} =
    -A \int_s^\infty
        x^{-2}
        \exp\left(-x^2 - \frac{2 \varepsilon \beta^3}{x}\right)
    \dd x.
\end{equation}

We obtain the constant of integration by the condition on the first derivative of $\dimless{T}$ at the interface:
\begin{equation}
    \dimless{T} =
    -(\xi + \omega \nu \tau) 2 \beta^3 \exp\left(\beta^2 + 2\varepsilon\beta^2\right)
    \int_s^\infty
        x^{-2}
        \exp\left(-x^2 - \frac{2 \varepsilon \beta^3}{x}\right)
    \dd x
\end{equation}

Finally, we have the condition that
\begin{equation}
    \begin{split}
        \text{Ja} &= \frac{\tau}{\xi + \omega \nu \tau} = \phi(\beta, \beta) \\
        \phi(a, b) &= g(a) f(a, b) \\
        g(a) &= 2 a^3 \exp\left((1+2\varepsilon) a^2\right) \\
        f(a, b) &= \int_b^\infty
            x^{-2}
            \exp\left(-x^2 - \frac{2 \varepsilon a^3}{x}\right)
        \dd x
    \end{split}
\end{equation}
from which we can calculate $\beta$.

\begin{equation}
    \text{Ja}
    = \frac{\rho_l c_l (T_\infty - T_\text{sat})}
           {\rho_g \left[L + (c_l - c_g) (T_\infty - T_\text{sat})\right]}
\end{equation}

\section{Two-Dimensional Scriven Problem}

\subsection{Conservation of Mass}

The two-dimensional rotationally symmetric continuity equation is
\begin{equation}
    \begin{split}
        &\frac{1}{r} \pd{r u_r}{r} + \frac{1}{r} \cancelto{0}{\pd{u_\theta}{\theta}} = 0 \\
        \Rightarrow\ & \frac{1}{r} \pd{r u_r}{r} = 0 \\
        \Rightarrow\ & \pd{u_r}{r} + \frac{u_r}{r} = 0
    \end{split}
\end{equation}
From this we conclude $ur = f(t)$, again dropping the $r$ subscript.

Taking the two-dimensional mass balance at the interface
\begin{equation}
    \begin{split}
        &\pd{}{t} \left(\pi R^2 \rho_g\right) = 
        2\pi R \rho_l \left[\dot{R} - u(R)\right] \\
        \Rightarrow\ & 2\pi R \dot{R} \rho_g = 
        2\pi R \rho_l \left[\dot{R} - u(R)\right] \\
        \Rightarrow\ & u(R) = 
        \frac{\rho_l - \rho_g}{\rho_l} \dot{R} = \varepsilon \dot{R}
    \end{split}
\end{equation}
we get the exact same result for the liquid velocity at the interface as for the three-dimensional case.

We again use this to evaluate the velocity in the entire domain.
\begin{equation}
    ur = \varepsilon \dot{R} R
    \label{eq:velocity_condition_2d}
\end{equation}
The main difference in the two-dimensional case is the missing square in the radius term.

\subsection{Conservation of Energy}

In the two-dimensional case, the rotationally symmetric conservation of energy is
\begin{equation}
    \pd{T}{t}
    =
    \alpha_l \left(
        \pdd{T}{r} +
        \frac1r \pd{T}{r}
    \right)
    - u \pd{T}{r}
    + \frac{Q}{\rho_l c_l}
\end{equation}

Using Eq. \eqref{eq:velocity_condition_2d}, we get
\begin{equation}
    \pd{T}{t}
    =
    \alpha_l \left(
        \pdd{T}{r} +
        \frac1r \pd{T}{r}
    \right)
    - \frac{\varepsilon \dot{R} R}{r} \pd{T}{r}
    + \frac{Q}{\rho_l c_l}
\end{equation}
which differs by a factor $2$ in the diffusion and a square in the velocity.

\subsection{Similarity Solution}

We apply the same simplifications as for the three-dimensional case and write the temperature equation in dimensionless form.
\begin{equation}
    \pd{\dimless{T}}{t} =
    \alpha_l \left(
        \pdd{\dimless{T}}{r} + \frac1r \pd{\dimless{T}}{r}
    \right)
    - \frac{\varepsilon \dot{R} R}{r} \pd{\dimless{T}}{r}
\end{equation}
We use the same definitions in Eqs.
\eqref{eq:dimensionless_superheat},
\eqref{eq:dimensionless_latent_heat},
\eqref{eq:dimensionless_diff_heat_capacity},
\eqref{eq:ratio_density},
\eqref{eq:similarity_r}, and
\eqref{eq:similarity_s}.

The we get
\begin{equation}
    \pdd{\dimless{T}}{s}
    =
    -\pd{\dimless{T}}{s} \left(
        2 s 
        + \frac{1 - 2 \varepsilon \beta^2}{s}
    \right)
\end{equation}
which has a factor $\frac12$ for the second term in the brackets, a square instead of a cube in the last term, and dividing by $s$ instead of $s^2$.

We again integrate over $s$ and get
\begin{equation}
    \pd{\dimless{T}}{s} = A s^{2\varepsilon \beta^2 - 1} \exp\left(-s^2\right)
\end{equation}
Integrating this equation again gives us
\begin{equation}
    \begin{split}
        \dimless{T}
        &= -A \int_s^\infty s^{2\varepsilon\beta^2 - 1} \exp\left(-s^2\right) \dd s \\
        &= -\frac{A}{2} \int_{s^2}^\infty x^{\varepsilon\beta^2 - 1} \exp\left(-x\right) \dd x \\
        &= -\frac{A}{2} \Gamma(\varepsilon \beta^2, s^2)
    \end{split}
\end{equation}
where $\Gamma$ is the incomplete upper gamma function.
We can do this because $\dimless{T}(\infty) = 0$.

On the interface we have the boundary condition
\begin{equation}
    \pd{\dimless{T}}{s}\bigg|_\beta = 2 \beta (\xi + \omega \nu \tau)
\end{equation}
from which we get
\begin{equation}
    A = 2 (\xi + \omega \nu \tau)
        \beta^{2 - 2 \varepsilon \beta^2}
        \exp\left(\beta^2\right).
\end{equation}
and finally
\begin{equation}
    \dimless{T}(s) =
    -(\xi + \omega \nu \tau)
        \beta^{2 - 2 \varepsilon \beta^2}
        \exp\left(\beta^2\right)
        \Gamma(\varepsilon \beta^2, s^2)
\end{equation}

Additionally, we have the condition at the interface
\begin{equation}
    \begin{split}
        &\dimless{T}(\beta) = -\tau \\
        \Rightarrow\ &
        \beta^{2 - 2 \varepsilon \beta^2}
          \exp\left(\beta^2\right)
          \Gamma(\varepsilon \beta^2, \beta^2)
        = \frac{\tau}{\xi + \omega \nu \tau} = \text{Ja}
    \end{split}
\end{equation}

To evaluate the temperature, we can make the following observation:
\begin{equation}
    \begin{split}
        &\dimless{T}(\beta) = -\tau \\
        \Rightarrow\ &
        (\xi + \omega \nu \tau)
        \beta^{2 - 2 \varepsilon \beta^2}
        \exp\left(\beta^2\right)
        = \frac{\tau}{\Gamma(\varepsilon \beta^2, \beta^2)} \\
        \Rightarrow\ &
        \dimless{T}(s) =
        -\tau \frac{\Gamma(\varepsilon \beta^2, s^2)}{\Gamma(\varepsilon \beta^2, \beta^2)} \\
        \Rightarrow\ &
        \dimless{T}(s) =
        -\tau \frac{\Gamma(\varepsilon \beta^2, s^2)}
                   {\Gamma(\varepsilon \beta^2, \beta^2)} =
        \frac{T_\text{sat} - T_\infty}{T_\infty}
        \frac{\Gamma(\varepsilon \beta^2, s^2)}
             {\Gamma(\varepsilon \beta^2, \beta^2)} \\
        \Rightarrow\ &
        T(s) =
        (T_\text{sat} - T_\infty)
        \frac{\Gamma(\varepsilon \beta^2, s^2)}
             {\Gamma(\varepsilon \beta^2, \beta^2)} + T_\infty
    \end{split}
\end{equation}

\newpage
\printbibliography

\end{document}